% =============================================================================
% preamble.tex — การตั้งค่าส่วนหัวสำหรับ "รายงานวิจัย 5 บท" ตามมาตรฐานวิชาการไทย
% คอมไพล์ด้วย XeLaTeX หรือ LuaLaTeX เท่านั้น (ต้องใช้ fontspec + polyglossia)
% =============================================================================

% ---------- ฟอนต์ภาษาไทย ----------
\usepackage{fontspec}
\usepackage{polyglossia}
\setdefaultlanguage{thai}
\setotherlanguage{english}

% เปิดใช้การตัดคำภาษาไทยแบบพจนานุกรมของ XeTeX (ICU) — จำเป็นมาก!
% ถ้าไม่เปิด บรรทัดข้อความไทยที่ไม่มีช่องว่าง (พบได้ทั่วไป โดยเฉพาะในช่องตาราง
% แคบ ๆ) จะไม่ตัดคำเลยและล้นออกนอกกรอบ/คอลัมน์แทนที่จะขึ้นบรรทัดใหม่
\XeTeXlinebreaklocale "th"
\XeTeXlinebreakskip = 0pt plus 1pt minus 0.1pt

% ฟอนต์หลัก TH Sarabun New ตามมาตรฐานราชการไทย (มติ ครม.)
% หากใช้ Overleaf: อัปโหลดไฟล์ฟอนต์ THSarabunNew.ttf / THSarabunNew Bold.ttf ฯลฯ
% ไปไว้ในโฟลเดอร์ fonts/ ของโปรเจกต์ (ดาวน์โหลดฟรีจาก SIPA/ f0nt.com)
% หากไม่มีไฟล์ฟอนต์ ให้สลับไปใช้ฟอนต์สำรอง "Sarabun" (มีอยู่แล้วบน Overleaf/Google Fonts)
% โดยคอมเมนต์บรรทัด TH Sarabun New แล้วเปิดใช้งานบรรทัด Sarabun แทน
\setmainfont{TH Sarabun New}[
  Path = fonts/,
  Extension = .ttf,
  UprightFont = *,
  BoldFont = * Bold,
  ItalicFont = * Italic,
  BoldItalicFont = * Bold Italic,
]
% ทางเลือกสำรอง (ฟอนต์ Sarabun จาก Google Fonts มีอยู่แล้วใน Overleaf บางเทมเพลต):
% \setmainfont{Sarabun}

\newfontfamily\thsarabun{TH Sarabun New}[Path = fonts/, Extension = .ttf,
  UprightFont = *, BoldFont = * Bold, ItalicFont = * Italic, BoldItalicFont = * Bold Italic]

% ---------- ขนาดหน้ากระดาษ/ระยะขอบ (ตรงตามไฟล์ตัวอย่างที่แนบ) ----------
\usepackage[a4paper,
  top=3.00cm, bottom=2.54cm, left=3.81cm, right=2.54cm,
  headsep=1.25cm, footskip=1.25cm]{geometry}

% ---------- ขนาดตัวอักษรมาตรฐาน ----------
% Normal = 16pt, หัวข้อรอง = 18pt ตัวหนา, บทที่ = 20pt ตัวหนา
\newcommand{\bodysize}{\fontsize{16}{24}\selectfont}
\newcommand{\hThreeSize}{\fontsize{16}{24}\selectfont\bfseries}
\newcommand{\hTwoSize}{\fontsize{18}{27}\selectfont\bfseries}
\newcommand{\hOneSize}{\fontsize{20}{30}\selectfont\bfseries}
\AtBeginDocument{\bodysize}

% ---------- ระยะบรรทัด/ย่อหน้า ----------
\usepackage{setspace}
\setstretch{1.25}
\usepackage{indentfirst}
\setlength{\parindent}{1.25cm}

% ---------- หัวข้อบท/หัวข้อย่อย (บทที่ N ให้ขึ้นหน้าใหม่ กึ่งกลาง ตัวหนา 20pt) ----------
\usepackage{titlesec}
\renewcommand{\thechapter}{\arabic{chapter}}
\titleformat{\chapter}[display]
  {\normalfont\hOneSize\centering}
  {บทที่ \thechapter}
  {8pt}
  {\hOneSize\centering}
\titlespacing*{\chapter}{0pt}{0pt}{24pt}

\titleformat{\section}
  {\normalfont\hTwoSize}{\thesection}{0.6em}{}
\titlespacing*{\section}{0pt}{18pt}{6pt}

\titleformat{\subsection}
  {\normalfont\hThreeSize}{\thesubsection}{0.6em}{}
\titlespacing*{\subsection}{0pt}{12pt}{6pt}

% ---------- สารบัญ / คำนำหน้าเอกสาร ----------
\usepackage{tocloft}
\renewcommand{\cftchapfont}{\bodysize}
\renewcommand{\cftsecfont}{\bodysize}
% ใช้ \AtBeginDocument เพื่อให้ชื่อหัวข้อไทยเหล่านี้ "ทับ" ค่าเริ่มต้นของ polyglossia
% (polyglossia ตั้งชื่อ \listfigurename ฯลฯ เป็นภาษาไทยของตัวเองตอนเริ่มเอกสารเช่นกัน
% ถ้าตั้งตรง ๆ ในพรีแอมเบิลเฉย ๆ จะถูก polyglossia ทับกลับ)
\AtBeginDocument{%
  \renewcommand{\contentsname}{สารบัญ}%
  \renewcommand{\listtablename}{สารบัญตาราง}%
  \renewcommand{\listfigurename}{สารบัญภาพ}%
  \renewcommand{\bibname}{บรรณานุกรม}%
  \renewcommand{\appendixname}{ภาคผนวก}%
}

% ---------- ตาราง/รูปภาพ ----------
\usepackage{longtable}
\usepackage{array}
\usepackage[table]{xcolor}
\usepackage{graphicx}
\usepackage{caption}
% หมายเหตุ: ใช้คีย์เวิร์ดฟอนต์มาตรฐานของแพ็กเกจ caption (ไม่ใช่มาโครที่มีอาร์กิวเมนต์
% อย่าง \fontsize{16}{24}) เพราะเมื่อใช้ร่วมกับ longtable (ltcaption) มาโครที่รับ
% อาร์กิวเมนต์จะทำให้ xelatex ค้าง/รายงาน "Incomplete \iffalse" ได้
\captionsetup{labelsep=quad, font=normalsize, labelfont=bf}
\AtBeginDocument{%
  \renewcommand{\figurename}{ภาพที่}%
  \renewcommand{\tablename}{ตารางที่}%
}

\definecolor{tblheadgray}{gray}{0.85}
\newcommand{\tblhead}[1]{\cellcolor{tblheadgray}\textbf{#1}}

% ---------- หัว-ท้ายกระดาษ / เลขหน้า ----------
\usepackage{fancyhdr}
\pagestyle{fancy}
\fancyhf{}
\fancyfoot[C]{\bodysize\thepage}
\renewcommand{\headrulewidth}{0pt}
\fancypagestyle{plain}{%
  \fancyhf{}
  \fancyfoot[C]{\bodysize\thepage}
  \renewcommand{\headrulewidth}{0pt}
}

% ---------- ระบบเลขหน้าพยัญชนะไทย (ก ข ค ...) สำหรับส่วนหน้า ----------
% ใช้ค่ารูปแบบ "thaiLetters" แบบเดียวกับที่ Microsoft Word ใช้กับส่วนหน้าของวิทยานิพนธ์ไทย
\usepackage{expl3}
\ExplSyntaxOn
\seq_new:N \g_thaifront_letters_seq
\seq_gset_split:Nnn \g_thaifront_letters_seq { , }
  {ก,ข,ค,ง,จ,ฉ,ช,ซ,ฌ,ญ,ฎ,ฏ,ฐ,ฑ,ฒ,ณ,ด,ต,ถ,ท,ธ,น,บ,ป,ผ,ฝ,พ,ฟ,ภ,ม,ย,ร,ล,ว,ศ,ษ,ส,ห,ฬ,อ,ฮ}
\NewDocumentCommand{\ThaiFrontPage}{m}{\seq_item:Nn \g_thaifront_letters_seq {#1}}
\ExplSyntaxOff

% เรียกใช้หลังหน้าปก เพื่อเริ่มเลขหน้า ก ข ค ... (นับใหม่จาก 1)
\newcommand{\thaifrontmatter}{%
  \cleardoublepage
  \pagenumbering{arabic}
  \renewcommand{\thepage}{\ThaiFrontPage{\arabic{page}}}%
}
% เรียกใช้ก่อนบทที่ 1 เพื่อเปลี่ยนเป็นเลขอารบิก 1 2 3 ... (นับใหม่จาก 1)
\newcommand{\mainbodymatter}{%
  \cleardoublepage
  \renewcommand{\thepage}{\arabic{page}}%
  \pagenumbering{arabic}%
}

% ---------- อ้างอิง ----------
\usepackage[style=apa, backend=biber]{biblatex}
\addbibresource{references.bib}

\usepackage{hyperref}
\hypersetup{hidelinks, bookmarksopen=true, pdfstartview=FitH}
